Autonomous navigation for mobile agricultural robots in horticultural environments faces severe operational challenges under Global Navigation Satellite Systems (GNSS). Thick plant canopies, multi-path reflections, and greenhouse structures degrade or eliminate satellite reception, while Real-Time Kinematic (RTK) corrections incur high infrastructure costs. High-resolution 3D LiDAR sensors offer direct geometric measurements, but unit costs remain prohibitive for scalable commercial fleets, and spatial point clouds lack the chromatic and textural cues required to distinguish drivable soil from dense crop foliage.

This thesis presents a real-time, deterministic vision-based navigation system for row-following horticulture based on monocular RGB imagery. The architecture decouples offline model-assisted dataset curation from low-latency onboard execution. To minimize manual annotation overhead, the Segment Anything Model (SAM~3) generates initial binary corridor segmentations from prompt-guided field captures. For real-time inference on an embedded NVIDIA Jetson Orin Nano platform, a lightweight MobileNetV3 + LR-ASPP semantic segmentation network is deployed via ONNX Runtime with FP16 precision, executing at 20~Hz within a 15~ms inference budget.

From the predicted drivable corridor mask, navigation setpoints are extracted using a decoupled dual-edge geometric tracking algorithm. The region of interest is partitioned into horizontal scanline strips from bottom to top. Left and right physical crop boundaries are tracked independently using spatial continuity gates and fitted using Iteratively Reweighted Least Squares (IRLS) with a Cauchy M-estimator. The central navigation line is computed analytically from the converged boundary parameters.

The proposed algorithm is benchmarked against three established agricultural vision baselines: a global column-projection extractor (SegMin column histogram), a strip-centroid regression line extractor, and a contour-point line extractor fitted via Random Sample Consensus (RANSAC). Rather than claiming unconditional superiority, the thesis evaluates the structural trade-offs of each paradigm. The proposed dual-edge extractor eliminates the stochastic frame-to-frame angular jitter characteristic of RANSAC ($\sigma_\omega = 0.00^\circ$ versus non-zero variance) and prevents unilateral leaf intrusions from biasing the corridor axis, a primary failure mode of centroid-based fitting. Conversely, experimental failure analysis formalizes key structural limitations of the proposed method: vulnerability to near-field seed occlusion, sensitivity of image-space jump gates to abrupt vehicle pitch and roll, inability of the straight-line model to track curved crop geometries or headland turnaround maneuvers, and boundary loss during extended unilateral crop gaps.

The perception and geometric extraction stages are validated offline across recorded horticultural datasets and benchmarked on the embedded NVIDIA Jetson Orin Nano, confirming real-time execution at 20~Hz with sub-15~ms latency without requiring secondary neural networks, stereo disparity maps, or LiDAR sensors. Finally, the thesis establishes the integrated control architecture and closed-loop field evaluation protocol on the Demetra four-wheel-steering (4WS) agricultural rover, using an auxiliary laser sensor to monitor physical corridor centering.